\chapter{Introduction}

\section{Motivation}

\subsection{Do humans dream of electric sheep?}

The dream of creating artificial beings that perceive and act like us is far older than the invention of the semiconductor.
From the ancient Greek myth of the bronze automaton Talos~\cite{mayor2018gods},
through the mechanical contraptions of the Enlightenment~\cite{riskin2003defecating,voskuhl2013androids},
to the very word \emph{robot}, coined in Karel Čapek's 1920 play R.U.R.~\cite{capek1920rur},
and the \emph{Maschinenmensch} (machine-human) of Fritz Lang's 1927 movie Metropolis~\cite{lang1927metropolis},
the idea of building mechanical creatures has been a recurring theme in our history.
Following the development of modern electronics and digital computers, this dream seemed within reach.
Throughout the mid-20th century, Isaac Asimov's tales of positronic-brain-powered robots~\cite{asimov1950irobot}
presented a vision of highly intelligent anthropomorphic machines that reason based on logical rules and hardcoded ethics to guide their actions. 
Stanisław Lem's stories~\cite{lem1965cyberiada} went a step further, giving them personalities, flaws, and desires,
showcasing the complexity of artificial civilisations and their potential ethical dilemmas.
Movies such as Forbidden Planet (1956)~\cite{wilcox1956forbidden} and Star Wars (1977)~\cite{lucas1977starwars}
presented robots as beings with an almost human psychology and thinking capabilities, albeit clumsier and more limited in their actions, bringing the idea to the mainstream.
On a philosophical level, the discourse shifted from Descartes' reflections that the
human living body might be nothing but an intricate machine~\cite{descartes1664homme} to Turing's question of whether we can actually determine if a machine is intelligent or not~\cite{turing1950computing}.
At the time, the founders of artificial intelligence confidently predicted that
human-level perception and reasoning were only a few years away~\cite{mccarthy1955proposal}.

Reality proved to be quite different and more humbling. Tasks that were regarded as the pinnacle of human intellect, such as playing chess~\cite{campbell2002deep} 
or proving mathematical theorems~\cite{davis1960computing,de2008z3}, turned out to be relatively easy for machines to master at a level surpassing even the best human experts.
At the same time, everyday competences such as finding a path in a cluttered room, locating and moving objects from one place to another, or even just looking for a familiar face in a crowd, turned out to be surprisingly
hard to reproduce even at the animal level. This inversion, known as Moravec's paradox~\cite{moravec1988mind}, has a compelling explanation in evolutionary
terms. Those abilities are a product of hundreds of millions of years of natural selection. The sudden emergence of eyes is even
credited with igniting the Cambrian explosion of animal life~\cite{parker2003blink,land2012animal}, and a large fraction of the
primate cerebral cortex is dedicated to processing what we see~\cite{felleman1991distributed}. We are, in effect, expert visual explorers by nature. What our brains do
almost subconsciously -- allocating gaze and gathering just enough visual evidence
to act -- remains as of today one of the hardest open problems in artificial intelligence~\cite{dreyfus1992computers}.

\subsection{The narrow window of perception}

Every agent that acts in the physical world, whether it is a human, an animal, or a machine, does so through a narrow window.
An eagle circling above a meadow searching for prey, a person looking for a familiar face in a crowded corridor, or a robot inspecting a factory floor for defects,
never perceives its surroundings all at once and in perfect detail. Instead, it is limited by its perceptual capabilities and environmental constraints. Furthermore,
the environment is not static, but changes constantly, making it impossible to map and analyse it in its entirety.
Therefore, the agent must be able to reason based on partial information, while actively pointing its senses where they matter to gather more data as needed.
The perception process is not a passive act of reception, but rather an active, goal-driven process of exploration~\cite{bajcsy1988active,aloimonos1988active,ballard1991animate}.
Human studies have shown that this exploration process happens on both a local and a global scale. 
On the local scale, we are experts at scanning our surroundings with our eyes, moving our gaze in a task-driven manner~\cite{yarbus1967eye,hayhoe2005eye}.
On the global scale, we excel at exploring our environment by moving our body to different locations, navigating through complex scenes and avoiding obstacles without much effort~\cite{fajen2003behavioral}.
The main motivation of this thesis is to enable machines to do the same: to \emph{learn to actively explore} their environment efficiently, perceiving only what they need in order to act on their goals.

The gap between this goal and the current state of machine learning and computer vision remains far from closed. Modern models achieve remarkable
accuracy in passive perception tasks~\cite{krizhevsky2012imagenet,dosovitskiy2021image,achiam2023gpt},
yet the very assumption of complete access to visual information~\cite{girdhar2022omnivore} breaks down the moment
a model is placed inside an embodied agent with a limited field of view, a
finite time budget, and a world that refuses to wait~\cite{ramakrishnan2021exploration,wenzel2021four}. Understanding how an agent
should efficiently acquire visual information, rather than merely how it should process information it has already
captured, remains an open problem~\cite{ramakrishnan2021exploration,bajcsy2018revisiting,duan2022survey,liu2025embodied}. The remainder of this chapter traces how the field arrived at this point, explains why general-purpose exploration is
still unsolved, and introduces the research direction pursued across the four publications that constitute this thesis.

\section{A brief outline of the field}

\subsection{From hand-crafted features to foundation models}

Computer vision is almost as old as computing itself, with the first works on machine perception dating back to the 1960s~\cite{roberts1963machine},
epitomised by the so-called ``blocks world'' of simple polyhedral scenes. First attempts to apply machine vision to real-world robotics problems followed
soon after, with the SRI Shakey project~\cite{nilsson1984shakey}, developed already in the late 1960s and early 1970s, and, roughly a decade later, with
Moravec's stereo-vision-guided Stanford Cart~\cite{moravec1980obstacle}. These early systems were built on hand-crafted geometry and operated in carefully prepared evaluation environments.
Their brittleness outside such controlled settings exposed how hard vision really is. Marr's influential computational theory of vision
\cite{marr1982vision} reframed the problem as one of recovering explicit representations of the world through successive stages of processing, and set an
agenda that would occupy the field for decades.

For the next two decades, progress was incremental at best and heavily reliant on hand-crafted feature representations.
The turning point came with the resurgence of neural networks, driven by the
availability of large annotated datasets~\cite{deng2009imagenet} and the development of more powerful hardware. Convolutional networks, pioneered for
digit recognition~\cite{lecun1998gradient}, were scaled up and decisively outperformed hand-crafted pipelines~\cite{krizhevsky2012imagenet}. Deeper and more trainable
architectures such as residual networks~\cite{he2016deep} followed, and learned representations rapidly became the default across tasks such as classification, detection, and
segmentation~\cite{lecun2015deep}. Computer vision had become, first and foremost, a problem of representation learning.

The next wave of progress came from the field of natural language processing, with the invention of attention-based transformer models~\cite{vaswani2017attention}.
The transformer architecture proved equally powerful for visual data, quickly surpassing convolutional networks in accuracy~\cite{dosovitskiy2021image}.
Moreover, transformers were able to learn rich feature representations with self-supervised and weakly-supervised learning~\cite{radford2021learning,caron2021emerging,he2022masked},
enabling training of models on vast amounts of unlabelled or loosely labelled data from the web to produce models with broad, general-purpose visual competence.
Those so-called vision foundation models in turn not only proved to be highly accurate backbones for standard vision tasks,
but also enabled the development of vision-language models~\cite{achiam2023gpt,team2023gemini}, which generalised and scaled up the ability of machines to transcribe images into text and vice versa to an unprecedented degree.

The trajectory from the 1960s blocks world to today's foundation models is, by any measure, a triumph. Yet nearly all of these
advances share a quiet assumption: that the model is handed all needed visual information in advance, in a complete, static, sufficiently high-resolution image.

\subsection{The passive-perception assumption}

Standard computer vision solutions assume complete and unobstructed access to input data~\cite{girdhar2022omnivore}. This is an entirely reasonable simplification for most studied tasks,
such as image classification or segmentation, and even the most recent vision-language benchmarks~\cite{chen2024we,yue2025mmmu} follow this pattern.
However, this assumption falls apart when the model is placed inside an \emph{embodied} agent, such as a robot or an unmanned aerial vehicle (UAV), operating in the real world.
Such an agent perceives the world through sensors with a restricted field of view and resolution, must move physically to acquire new observations, taking obstacles into account, and
operates under tight constraints on time, energy, and computation power~\cite{wenzel2021four}. Capturing and processing every part of a large environment at maximum resolution is not merely wasteful but often infeasible, especially in open-world settings~\cite{mnih2014recurrent,sandini2003retina,ramakrishnan2021exploration}.

Under these conditions, the question is no longer only \emph{what we currently see}, but rather \emph{what we should observe} in the first
place. An agent that can choose its observations wisely can achieve strong performance at a fraction of the sensing and computational cost of one that
processes everything indiscriminately. For active perception to work we need to move this decision of \emph{where to look next} inside the model itself.

\subsection{Active embodied vision}

The idea that perception systems should be active has been around since the early days of embodied computer vision.
The active perception~\cite{bajcsy1988active}, active vision~\cite{aloimonos1988active}, and animate vision~\cite{ballard1991animate} proponents argued decades ago that a perceiver
which controls its own sensing can turn otherwise ill-posed problems into tractable ones, and this perspective has been revisited in the modern learning
era \cite{bajcsy2018revisiting}. The inspiration is deeply biological. Human vision is highly foveated, concentrating acuity at the centre of gaze. Therefore we
explore scenes through rapid saccadic eye movements that direct this fovea to informative regions~\cite{yarbus1967eye,hayhoe2005eye}.

Machine learning approaches to active vision followed the broader deep learning progress.
Recurrent models of visual attention~\cite{mnih2014recurrent} learned to select a sequence of glimpses for recognition,
drawing on the reinforcement learning methods~\cite{sutton2018reinforcement} that had begun to master control
problems from raw pixels~\cite{mnih2015human}. In parallel, the robotics and embodied-AI communities studied exploration in the context of simultaneous
localisation and mapping (SLAM)~\cite{cadena2016past}, where an agent must build an explicit map of the environment.
More recently, the task of embodied visual exploration~\cite{ramakrishnan2021exploration} studied how an agent should build an internal model of an unfamiliar space to support downstream tasks such as reconstruction, localisation, and navigation.
 Another recent line of work on active visual exploration specifically studies how a model equipped with a
limited field of view should select successive glimpses to reconstruct,
classify, or segment a scene \cite{seifi2021glimpse}. These threads share a common conviction: that intelligent sensing is a decision problem.

\subsection{Why general exploration is still unsolved}

Despite this long history of research and recent rapid progress, the problem of exploration in the general case remains far from solved.
The difficulty is easy to underestimate because existing methods achieve strong performance within simplified settings.
For example, some active-vision methods operate in small search spaces, where even a random strategy can achieve strong performance~\cite{seifi2021glimpse}.
Moreover, the exploration is often limited to a grid game, where an agent uncovers a fixed area at each move~\cite{rangrej2022consistency},
 a simplification that mirrors the grid input format of vision transformers rather than the continuous movements of real-world agents. Embodied navigation is most often studied in constrained indoor simulators~\cite{savva2019habitat,kolve2017ai2,xia2018gibson},
where the agent can map the entire environment explicitly, without having to deal with the complexity and vast size of real-world environments.
This allows relatively simple object-goal navigation methods~\cite{batra2020objectnav,gu2022vision} to produce strong results on those benchmarks, but it does not address the problem of general exploration.
A different body of work pushes this global-scale exploration outdoors, using simulators such as AirSim~\cite{airsim2017fsr}, vision-and-language navigation tasks such as
AerialVLN~\cite{liu2023aerialvln} and CityNav~\cite{lee2024citynav}, or, more rarely, open-ended object-goal search directly in such settings~\cite{xie2023reasoning}.
Notably, most of this active-vision progress, indoors and outdoors alike, still operates at what we earlier called the
local scale, or follows a fairly constrained navigation task: an agent at each step decides where to look within a scene it can fully scan, or move towards an object it can already see.
Whether the same efficient, learned decision-making transfers to unconstrained, zero-shot, global-scale search remains largely an open question~\cite{bajcsy2018revisiting,ramakrishnan2021exploration}.
Finally, it is not clear whether methods designed for simplified simulation environments can be transferred to the real physical world, bridging the sim-to-real gap~\cite{aljalbout2025reality}.


\section{Learning to explore efficiently}

The central thesis of this dissertation is that \emph{active visual exploration} (AVE) should be treated as a learning problem in its own right.
Doing so yields agents which, first of all, are \emph{efficient}, in the sense of understanding an environment from as few observations as possible.
Second, they are \emph{flexible}, in the sense of exploiting the true continuous capabilities of real sensors rather than artificial grids.
Finally, they are \emph{scalable} toward the demands of real embodied agents operating in the open world, spanning both the local scale of choosing where to look and the global scale of choosing where to move.
The four publications that make up this thesis form a deliberate progression along this path, each dismantling a limitation left standing by the previous one.

The first contribution, ``Active Visual Exploration Based on Attention-Map Entropy''~\cite{pardyl2023active}, asks how an agent should decide where to look
next without the machinery of a separate decision network. It shows that the internal uncertainty of a transformer, read directly from the entropy of its
attention maps, is itself an effective signal for selecting the most informative next glimpse, simplifying training while improving reconstruction,
classification, and segmentation from partial observations. This work operates, however, within the familiar setting of glimpses drawn from a fixed grid.

The second contribution, ``Beyond Grids: Exploring Elastic Input Sampling for Vision Transformers''~\cite{pardyl2023grids}, addresses that constraint.
It formalises the notion of input \emph{elasticity}, the resilience of a model to patches of varying scale, position, and coverage, introduces a protocol to
measure it, and proposes architectural and training modifications that let a vision transformer ingest patches sampled freely from an image. This turns the rigid grid into a continuous sampling space.

The third contribution, ``AdaGlimpse: Active Visual Exploration with Arbitrary Glimpse Position and Scale''~\cite{pardyl2024adaglimpse}, closes the gap between
software and hardware. Building on an input-elastic transformer, it casts exploration as a Markov decision process and uses the Soft Actor-Critic
reinforcement learning algorithm to select glimpses of arbitrary continuous position and scale, exactly the degrees of freedom offered by optical zoom or by
changing a UAV's altitude. The resulting agent rapidly establishes a coarse awareness of a scene before zooming in on detail, matching or surpassing prior
methods while observing far fewer pixels.

The first three papers progressively generalise \emph{local-scale} exploration: deciding
where to look within a scene the agent can already sense, moving from a fixed grid (AME), to
elastic patch geometry (Beyond Grids), to arbitrary continuous zoom and pan (AdaGlimpse). The
fourth contribution, ``FlySearch: Exploring how vision-language models explore''~\cite{pardyl2025flysearch}, turns to the complementary \emph{global-scale} problem: deciding where
to move next in a large, open, three-dimensional world in order to find something. The environment
is no longer a bounded scene from which glimpses can be sampled, but an unbounded outdoor space
that the agent, a simulated UAV, must physically traverse to acquire each new observation. FlySearch introduces a photorealistic benchmark for goal-driven object search in such
environments and uses it to ask whether today's vision-language models can execute this kind of exploration. It finds that
state-of-the-art models fail to reliably solve even simple exploration tasks,
with the gap to human performance widening sharply as tasks demand longer,
consistent search strategies, and this gap persists even after fine-tuning. In
mapping out these failure modes, FlySearch both charts the frontier of the
problem and motivates further research in this field.

Taken together, these works trace an arc from deciding where to look within a single reachable
scene, through freeing the model from rigid sampling and matching real sensor geometry, to deciding
where to move through an open, three-dimensional world. They constitute concrete steps toward the
goal that gives this thesis its title: learning to explore, at both the local scale of gaze and
the global scale of locomotion, on the way to efficient embodied vision.

\section{Thesis outline}

The remainder of this thesis follows the progression outlined above.
Chapter~\ref{chap:uncertainty} studies exploration at the local scale and presents the first publication, ``Active Visual Exploration Based on Attention-Map Entropy'', which recasts glimpse selection as uncertainty reduction driven by the internal attention entropy of the task model itself.
Chapter~\ref{chap:adaptive} lifts the fixed-grid assumption underlying that work and covers two publications: ``Beyond Grids: Exploring Elastic Input Sampling for Vision Transformers'', which formalises input elasticity and introduces the ElasticViT architecture, and ``AdaGlimpse: Active Visual Exploration with Arbitrary Glimpse Position and Scale'', which trains a reinforcement learning agent to select glimpses of continuous position and scale on top of that backbone.
Chapter~\ref{chap:openworld} moves to the global scale and presents ``FlySearch: Exploring how vision-language models explore'', a benchmark for goal-driven object search in large outdoor environments.
Each of these chapters opens with its research scope and the preliminaries needed to keep the thesis self-contained, then summarises the corresponding publication and its contributions; the full texts of the publications are provided in the appendix.
Chapter~\ref{chapter:other} gives an overview of other research work and achievements from the author's doctoral studies, and Chapter~\ref{chap:conclusion} concludes the thesis with a synthesis of the results and a discussion of future research directions.